\documentclass[11pt,a4paper]{article}
\usepackage[utf8]{inputenc}
\usepackage[T1]{fontenc}
\usepackage[french]{babel}
\usepackage{lmodern}
\usepackage{microtype}
\usepackage{geometry}
\usepackage{amsmath,amssymb,mathtools}
\usepackage{booktabs}
\usepackage{xcolor}
\usepackage{hyperref}
\geometry{margin=2.4cm}
\hypersetup{colorlinks=true,linkcolor=blue!60!black,citecolor=blue!60!black,urlcolor=blue!60!black}
\setlength{\parindent}{0pt}
\setlength{\parskip}{0.6em}

\title{Tests observationnels publics du cadre GUP--UV/IR}
\author{Note de verification}
\date{22 mai 2026}

\newcommand{\rhoP}{\rho_{\rm P}}
\newcommand{\rhoDE}{\rho_{\rm DE}}
\newcommand{\ellP}{\ell_{\rm P}}
\newcommand{\EGUP}{E_{\rm GUP}}

\begin{document}
\maketitle

\begin{abstract}
Cette note teste les consequences observationnelles simples du cadre GUP--UV/IR
avec des donnees ou contraintes publiques : FIRAS pour le spectre du CMB,
Planck+BAO+SNe pour l'equation d'etat de l'energie noire, SPARC pour l'echelle
d'acceleration galactique, BBN/CMB pour la densite de radiation et etoiles a
neutrons pour les fermions denses. Le resultat principal est un no-go : la
valeur $\beta_0\sim10^{60}$ obtenue en identifiant directement le vide GUP a
$\rhoDE$ est exclue comme deformation universelle. La piste viable est une
regularisation UV par GUP combinee a une projection IR/holographique du secteur
gravitationnel.
\end{abstract}

\section{Parametre teste}

La regularisation GUP du vide donne
\begin{equation}
\rho_{\rm reg}=\frac{g\rhoP}{16\pi^2\beta_0^2}.
\end{equation}
L'identification directe $\rho_{\rm reg}=\rhoDE$ impose
\begin{equation}
\beta_0^{(\Lambda)}
=
\left(\frac{g}{6\pi\Omega_\Lambda}\right)^{1/2}
\frac{1}{H_0t_P}.
\end{equation}
Pour $g=1$ et les parametres Planck 2018,
\begin{equation}
\beta_0^{(\Lambda)}=2.36\times10^{60},
\qquad
\EGUP=\frac{E_P}{\sqrt{\beta_0}}=7.94~{\rm meV}.
\end{equation}

\section{FIRAS : spectre du CMB}

Si le meme $\beta_0$ s'applique aux photons, le spectre thermique devient
\begin{equation}
u(E,T)\propto
\frac{E^3}{(e^{E/k_BT}-1)(1+\beta E^2/c^2)^3}.
\end{equation}
Le script compare cette forme au meilleur corps noir ajuste sur la bande FIRAS
2--20 cm$^{-1}$. FIRAS limite les deviations de forme a environ $50$ ppm du pic
du CMB.

Le resultat numerique est
\begin{equation}
\frac{\mathrm{RMS}}{\mathrm{pic}}=2.81\times10^{-3},
\qquad
\beta_0^{\rm FIRAS}\lesssim3.72\times10^{58}.
\end{equation}
La valeur $\beta_0^{(\Lambda)}$ est donc trop grande par un facteur environ
$64$. Donc
\begin{equation}
\boxed{
\beta_0^{(\Lambda)}\text{ est exclu comme parametre photon universel.}
}
\end{equation}

\section{Equation d'etat de l'energie noire holographique}

Pour un modele holographique standard fonde sur l'horizon des evenements futur,
\begin{equation}
w_{\rm HDE}
=
-\frac13-\frac{2}{3c_{\rm hde}}\sqrt{\Omega_{\rm DE}}.
\end{equation}
Avec $\Omega_{\rm DE}=0.685$, on obtient
\begin{equation}
w(c_{\rm hde}=1)=-0.885.
\end{equation}
La contrainte publique Planck+BAO+SNe donne environ $w=-1.03\pm0.03$. Ainsi,
$c_{\rm hde}=1$ est defavorise, tandis qu'une valeur proche de
$c_{\rm hde}\simeq0.79$ reproduit mieux la contrainte constante-$w$.

\section{SPARC et acceleration cosmologique}

L'echelle de de Sitter est
\begin{equation}
a_{\rm dS}=c^2\sqrt{\frac{\Lambda}{3}},
\end{equation}
et
\begin{equation}
\frac{a_{\rm dS}}{2\pi}=8.63\times10^{-11}~{\rm m\,s^{-2}}.
\end{equation}
L'echelle empirique SPARC/MOND est de l'ordre
\begin{equation}
a_0\simeq1.2\times10^{-10}~{\rm m\,s^{-2}}.
\end{equation}
L'accord est d'ordre unite. C'est une coincidence UV/IR interessante, mais pas
une preuve du GUP.

\section{BBN et densite de radiation}

La correction de rayonnement est
\begin{equation}
\frac{\Delta u}{u}
\simeq
-\frac{40\pi^2}{7}\beta_0\left(\frac{k_BT}{E_P}\right)^2.
\end{equation}
A $k_BT\simeq1$ MeV, imposer conservativement
$|\Delta u/u|<0.1$ donne
\begin{equation}
\beta_0\lesssim 2.64\times10^{41}.
\end{equation}
Pour $\beta_0^{(\Lambda)}$, le script trouve
$|\Delta u/u|\simeq8.94\times10^{17}$. Cette valeur est exclue dans le secteur
radiation du jeune univers.

\section{Etoiles a neutrons}

Pour une impulsion de Fermi representative $p_Fc\simeq300$ MeV, la correction
est
\begin{equation}
\chi=\beta_0\left(\frac{p_Fc}{E_P}\right)^2.
\end{equation}
Exiger $\chi<0.1$ donne une borne d'ordre
\begin{equation}
\beta_0\lesssim1.66\times10^{38}.
\end{equation}
La valeur $\beta_0^{(\Lambda)}$ donnerait $\chi\simeq1.43\times10^{21}$ ; elle
est donc incompatible avec une deformation universelle des fermions denses.

\section{Conclusion}

Les tests publics donnent une conclusion claire :
\begin{equation}
\boxed{
\beta_0^{(\Lambda)}\sim10^{60}
\text{ ne peut pas etre un parametre universel.}
}
\end{equation}
La bonne lecture des travaux est donc
\begin{equation}
\boxed{
\text{GUP regularise l'UV, mais l'energie noire doit venir d'une projection IR/holographique.}
}
\end{equation}

\begin{thebibliography}{9}
\bibitem{FIRAS}
COBE/FIRAS data products, NASA LAMBDA,
\url{https://lambda.gsfc.nasa.gov/product/cobe/firas_products.html}.

\bibitem{Planck2018}
N. Aghanim et al. (Planck Collaboration),
``Planck 2018 results. VI. Cosmological parameters'',
\emph{Astronomy \& Astrophysics} \textbf{641}, A6 (2020),
arXiv:1807.06209.

\bibitem{PantheonPlus}
Pantheon+SH0ES Data Release,
\url{https://github.com/PantheonPlusSH0ES/DataRelease}.

\bibitem{SPARC}
SPARC galaxy rotation-curve database,
\url{https://zenodo.org/records/16284118}.

\bibitem{CKN}
A. Cohen, D. Kaplan and A. Nelson,
``Effective Field Theory, Black Holes, and the Cosmological Constant'',
\emph{Physical Review Letters} \textbf{82}, 4971 (1999).
\end{thebibliography}

\end{document}
